\subsection{SPACES OF MAPPINGS INTO A NORMED SPACE}

Consider, more particularly, the situation in which Y is a normed vector space over a non-discrete valued division ring K (Chapter IX, § 3, no. 3). Let us denote by \(||y||\) the norm of \(y \in Y\) . The set \(\mathcal{F}(X;Y) = Y^{\mathbf{X}}\) is then canonically endowed with a K-vector space structure. A mapping \(u: X \to Y\) is bounded if and only if the real-valued function \(x \to ||u(x)||\) is bounded in X. If \(u, v\) are bounded mappings of X into Y, it is clear that \(u + v\) and \(\lambda u (\lambda \in K)\) are bounded; in other words, \(\mathcal{B}(X;Y)\) is a vector subspace of \(\mathcal{F}(X;Y)\) . Moreover, \(||u|| = \sup_{x \in X} ||u(x)||\) is a norm on \(\mathcal{B}(X;Y)\) ; for it satisfies the triangle inequality and \(||u|| = 0\) implies \(u = 0\) , and for each \(\lambda \in K\) we have

\[
| | \lambda u | | = \sup _ {x \in \mathbf {X}} | | \lambda u (x) | | = \sup _ {x \in \mathbf {X}} | \lambda |. | | u (x) | | = | \lambda |. \sup _ {x \in \mathbf {X}} | | u (x) | | = | \lambda |. | | u | |.
\]

Moreover, it is immediately verified that the uniformity on \(\mathcal{B}(\mathbf{X};\mathbf{Y})\) defined by this norm is the uniformity of uniform convergence. Unless the contrary is expressly stated, whenever \(\mathcal{B}(\mathbf{X};\mathbf{Y})\) is considered as a normed space, it is the norm defined above which is in question.

PROPOSITION 3. If the normed space Y is complete, then every series \((u_n)\) of bounded mappings of X into Y which is absolutely convergent in the normed space \(\mathcal{B}\) (X; Y) (i.e.~which is such that \(\sum_{n=0}^{\infty}||u_n|| < +\infty\) ; cf.~Chapter IX, § 3, no. 6) is uniformly convergent in X.

For since \(\mathcal{B}(\mathbf{X};\mathbf{Y})\) is complete (no. 1, Proposition 2, Corollary 1), the result follows from Chapter IX, § 3, no. 6, Proposition 11 and the definition of a uniformly convergent series.

\[
\text { Remark. } \quad 1) \text {   If   } \sum_ {n = 0} ^ {\infty} \| u _ {n} \| <   + \infty , \text {   then   } \sum_ {n = 0} ^ {\infty} \| u _ {n} (x) \| \leqslant \sum_ {n = 0} ^ {\infty} \| u _ {n} \| <   + \infty
\]

for each \(x \in \mathbf{X}\) ; in other words, for each \(x \in \mathbf{X}\) the series with general term \(u_n(x)\) is absolutely convergent in the space \(\mathbf{Y}\) . The converse is false. To avoid all confusion we shall sometimes say that the series with general term \(u_n\) is normally convergent, meaning that the series with general term \(\| u_n \|\) is convergent. A series can be uniformly convergent in \(\mathbf{X}\) without being normally convergent; this is the case, for example, for the series \((u_n)\) in the space \(\mathcal{B}(\mathbf{R}, \mathbf{R})\) , defined as follows: \(u_n(x) = (1/n) \sin x\) if \(x \in [n\pi, (n + 1)\pi]\) , \(u_n(x) = 0\) otherwise.

When Y is a normed algebra (Chapter IX, § 3, no. 7) over a non-discrete valued field K, then \(\mathcal{B}(\mathbf{X};\mathbf{Y})\) is a K-algebra, and the norm \(||u||\) is compatible with the algebra structure, since

\[
\begin{array}{r l} | | u v | | = \sup _ {x \in \mathbf {X}} | | u (x) v (x) | | & \leqslant \sup _ {x \in \mathbf {X}} | | u (x) | |. | | v (x) | | \\ & \leqslant \sup _ {x \in \mathbf {X}} | | u (x) | |. \sup _ {x \in \mathbf {X}} | | v (x) | | = | | u | |. | | v | |. \end{array}
\]

Thus \(\mathcal{B}(\mathbf{X};\mathbf{Y})\) is now a normed algebra over \(\mathbf{K}\) .

PROPOSITION 4. Let \(\mathbf{X}_i\) ( \(1 \leqslant i \leqslant n\) ) and \(\mathbf{Y}\) be normed vector spaces over a non-discrete valued division ring \(\mathbf{K}\) , and let \(\mathbf{X} = \prod_{i=1}^{n} \mathbf{X}_i\) . Then the set of all multilinear mappings of \(\mathbf{X}\) into \(\mathbf{Y}\) is closed in the space \(\mathcal{F}_s(\mathbf{X}; \mathbf{Y})\) .

This set consists of all \(u \in \mathbf{F}(\mathbf{X};\mathbf{Y})\) which satisfy all the relations

\[
\begin{array}{r l} u (x _ {1}, \dots , x _ {i} ^ {\prime} + x _ {i} ^ {\prime \prime}, \dots , x _ {n}) & = u (x _ {1}, \dots , x _ {i} ^ {\prime}, \dots , x _ {n}) \\ & \quad + u (x _ {1}, \dots , x _ {i} ^ {\prime \prime}, \dots , x _ {n}), \\ u (x _ {1}, \dots , \lambda x _ {i}, \dots , x _ {n}) & = \lambda u (x _ {1}, \dots , x _ {i}, \dots , x _ {n}) \end{array}\tag{1}
\]

\(1 \leqslant i \leqslant n, x_i, x_i', x_i''\) arbitrary elements of \(\mathbf{X}_i\) , \(\lambda\) an arbitrary element of \(\mathbf{K}\) ); since both sides of the relations (1) are continuous functions of \(u\) on \(\mathcal{F}_s(\mathbf{X}; \mathbf{Y})\) ( \(\S\) 1, no. 2, Remark 6), the result follows (Chapter I, § 8, no. 1, Proposition 2).

PROPOSITION 5. Under the hypotheses of Proposition 4, the set \(\mathcal{L}(\mathbf{X}_1,\ldots ,\mathbf{X}_n;\mathbf{Y})\) of continuous multilinear mappings of \(\mathbf{X}\) into \(\mathbf{Y}\) is closed in \(\mathcal{F}(\mathbf{X};\mathbf{Y})\) with respect to the topology of bounded convergence; it is complete with respect to the uniformity of bounded convergence if \(\mathbf{Y}\) is complete.

For if \(\mathfrak{S}\) is the set of all bounded subsets of \(\mathbf{X},\mathcal{L}(\mathbf{X}_1,\ldots ,\mathbf{X}_n;\mathbf{Y})\) is the intersection of the set of all multilinear mappings of \(\mathbf{X}\) into \(\mathbf{Y}\) and the set \(\mathcal{B}_{\mathfrak{S}}(\mathbf{X};\mathbf{Y})\) (Chapter IX, § 3, no. 5, Theorem 1); the result thus follows from Proposition 4 and Proposition 2, Corollary 1.

For the remainder of this subsection, K denotes a non-discrete valued field.

Then \(\mathcal{L}(\mathbf{X}_1, \ldots, \mathbf{X}_n; \mathbf{Y})\) is a vector subspace of \(\mathcal{F}(\mathbf{X}; \mathbf{Y})\) . Let \(\mathbf{B}\) be the unit ball in \(\mathbf{X}\) , the set of all \((\boldsymbol{x}_i)_{1 \leqslant i \leqslant n}\) such that \(\sup_{1 \leqslant i \leqslant n} ||\boldsymbol{x}_i|| \leqslant 1\) . Then the mapping \(u \to u|\mathbf{B}\) of \(\mathcal{L}(\mathbf{X}_1, \ldots, \mathbf{X}_n; \mathbf{Y})\) into \(\mathcal{B}(\mathbf{B}; \mathbf{Y})\) is injective; moreover, the inverse image, under this mapping, of the uniformity of uniform convergence on \(\mathcal{B}(\mathbf{B}; \mathbf{Y})\) is the uniformity of bounded convergence on \(\mathcal{L}(\mathbf{X}_1, \ldots, \mathbf{X}_n; \mathbf{Y})\) . For every bounded subset of \(\mathbf{X}\) is contained in a set of the form \(\mu\mathbf{B}\) (for some \(\mu \in \mathbf{K}^*\) ), and if \(u\) is an element of \(\mathcal{L}(\mathbf{X}_1, \ldots, \mathbf{X}_n; \mathbf{Y})\) , to say that \(||u(z)|| \leqslant a\) for all \(z \in \mu\mathbf{B}\) is equivalent to saying that \(||u(z)|| \leqslant a/|\mu|^n\) for all \(z \in \mathbf{B}\) . It is easily verified that the number

\[
\left| \left| u \right| \right| = \sup _ {z \neq 0} \frac {\left| \left| u (z) \right| \right|}{\left| \left| z \right| \right|}
\]

is a norm on \(\mathcal{L}(\mathbf{X}_1, \ldots, \mathbf{X}_n; \mathbf{Y})\) and defines the uniformity of bounded convergence on this set, and clearly we have

\[
\left| \left| u \left(x _ {1}, \dots , x _ {n}\right) \right| \right| \leqslant \left| \left| u \right| \right|. \left| \left| x _ {1} \right| \right| \dots \left| \left| x _ {n} \right| \right|.\tag{2}
\]

Unless the contrary is expressly stated, whenever \(\mathcal{L}(\mathbf{X}_1,\ldots ,\mathbf{X}_n;\mathbf{Y})\) is considered as a normed space, it is the norm defined above which is in question.

PROPOSITION 6. The multilinear mapping

\[
(u, x _ {1}, \dots , x _ {n}) \rightarrow u (x _ {1}, \dots , x _ {n})
\]

of the normed space \(\mathcal{L}(\mathbf{X}_{1},\ldots,\mathbf{X}_{n};\mathbf{Y})\times\mathbf{X}_{1}\times\cdots\times\mathbf{X}_{n}\) into Y is continuous.

This is an immediate consequence of the inequality (2) (Chapter IX, § 3, no. 5, Theorem 1).

PROPOSITION 7. Let X, Y, Z be three normed spaces over K. The canonical mapping of the normed space \(\mathscr{L}(X, Y; Z)\) into the space of linear mappings of X into \(\mathscr{L}(Y; Z)\) which sends each \(u \in \mathscr{L}(X, Y; Z)\) to the mapping \(x \to u(x, .)\) is an isometry of \(\mathscr{L}(X; Y; Z)\) onto \(\mathscr{L}(X; \mathscr{L}(Y; Z))\) .

This follows immediately from the definitions and the relation

\[
\sup _ {| | \boldsymbol {x} | | \leqslant 1} \left(\sup _ {| | \boldsymbol {y} | | \leqslant 1} | | u (\boldsymbol {x}, \boldsymbol {y}) | |\right) = \sup _ {| | \boldsymbol {x} | | \leqslant 1, | | \boldsymbol {y} | | \leqslant 1} | | u (\boldsymbol {x}, \boldsymbol {y}) | |.
\]

PROPOSITION 8. Let X, Y, Z be three normed spaces over K. The bilinear mapping \((u, v) \to v \circ u\) of \(\mathcal{L}(X; Y) \times \mathcal{L}(Y; Z)\) into \(\mathcal{L}(X; Z)\) is continuous.

For if \(u \in \mathcal{L}(\mathbf{X};\mathbf{Y})\) and \(v \in \mathcal{L}(\mathbf{Y};\mathbf{Z})\) we have

\[
\left| \left| v \circ u \right| \right| \leqslant \left| \left| u \right| \right| \cdot \left| \left| v \right| \right|,\tag{3}
\]

since for all \(\pmb{x} \in \mathbf{X}\) we have \(||v(u(\pmb{x}))|| \leqslant ||v|| \cdot ||u(\pmb{x})|| \leqslant ||v|| \cdot ||u|| \cdot ||\pmb{x}||\) by reason of (2).

In particular, on the set \(\mathcal{L}(\mathbf{X})\) of continuous endomorphisms of a normed space \(\mathbf{X}\) over \(\mathbf{K}\) , the norm \(||u||\) is compatible with the K-algebra structure of \(\mathcal{L}(\mathbf{X})\) .

Remark. 2) The set \(\mathcal{L}(\mathbf{R}^{m};\mathbf{R}^{n})\) of linear (necessarily continuous) mappings of \(\mathbf{R}^{m}\) into \(\mathbf{R}^{n}\) can be identified with the set \(\mathbf{M}_{n,m}(\mathbf{R})\) of matrices with \(n\) rows and \(m\) columns with coefficients in \(\mathbf{R}\) and hence can be identified with \(\mathbf{R}^{mn}\) ; on \(\mathcal{L}(\mathbf{R}^{m};\mathbf{R}^{n})\) , the uniformity of bounded convergence (with respect to the Euclidean metric on \(\mathbf{R}^{m}\) ), of compact convergence, and of pointwise convergence are then identified with the additive uniformity on \(\mathbf{R}^{mn}\) . Take the norm of \(x = (x_{i})\in \mathbb{R}^{n}\) to be

\[
\| \boldsymbol {x} \| = \sup _ {i} | x _ {i} |,
\]

and let \((\pmb{e}_j)\) be the canonical basis of \(\mathbf{R}^m\) ; if \(u\) and \(v\) are two linear mappings of \(\mathbf{R}^m\) into \(\mathbf{R}^n\) such that \(\| u(\pmb{e}_j) - v(\pmb{e}_j)\| \leqslant \varepsilon\) for \(1 \leqslant j \leqslant m\) , then we have \(|\alpha_{ij} - \beta_{ij}| \leqslant \varepsilon\) for each pair \((i,j)\) \([U = (\alpha_{ij})\) and \(V = (\beta_{ij})\) being the matrices of \(u, v\) respectively{]}; and conversely, if these inequalities are satisfied, we have \(\| u(\pmb{x}) - v(\pmb{x})\| \leqslant ma\varepsilon\) for every point \(\pmb{x}\) of a cube of centre 0 and side \(a\) in \(\mathbf{R}^m\) .

